\subsection{Schemes}

\begin{exer}
	Let $X$ be an open subscheme of an affine scheme $Y$. The canonical
	morphism $X \to Y$ corresponds by the map $\rho$ of Proposition 25 to
	the restriction $\Sh_Y(Y) \to \Sh_X(X)$.
\end{exer}
\begin{proof}
	In the case where $X$ is a principal open subset, this is just Lemma
	3.7. Otherwise, suppose that $X = \bigcup D(f_i)$ for some family of
	elements $f_i \in \Sh_Y(Y)$. The inclusions induce a family of maps
	$\Sh_Y(Y) \to \Sh_X(D(f_i)) = \Sh_X(X)_{f_i}$ which clearly agree on the
	intersections. Gluing these gives the restriction.
\end{proof}

\begin{exer}
	Let $U = \Spec B$ be an affine open subscheme of $X = \Spec A$. Then the
	restriction $A \to B$ is flat.
\end{exer}
\begin{proof}
	Since the restriction corresponds to the open immersion $U
	\hookrightarrow X$, we see that its localization $A_x \to B_x$ is an
	isomorphism, which is flat, for every $x \in \Spec B$. 
\end{proof}

\begin{exer} \leavevmode
	\begin{enumerate}[label=(\alph*)]
		\item Let $F$ be a subset of an affine scheme $\Spec A$. The
			closure $\overline F$ of $F$ in $\Spec A$ is equal to
			$V(I)$, where $I = \bigcap_{\mathfrak p \in F}
			\mathfrak p$.
		\item Let $\varphi: A \to B$ be a ring homomorphism. Let $f:
			\Spec B \to \Spec A$ be the morphism of schemes
			associated to $\varphi$. Then $\overline{\Im f} =
			V(\Ker \varphi).$
	\end{enumerate}
\end{exer}
\begin{proof}
	We already have that $\overline F$ must already be of the form $V(I')$
	with $I' \subset I$. Since $V(I')$ must be the minimal such set, we have
	$I' = I$.\\

	For part (b) we assume that $A$ and $B$ are reduced, since that doesn't
	affect their prime spectra. Then the previous part implies that
	$\overline{\Im f} = V(\bigcap_{\mathfrak p \in \Spec A} \varphi^{-1}
	(\mathfrak p)) = V(\Ker \varphi)$.
\end{proof}
In the case where $B = A_{\mathfrak p}$ for some $\mathfrak p \in \Spec A$,
this is the set of ideals that are "large enough" to interact with
$\mathfrak p$. In geometric terms, they are the irreducible varieties such that
every function that vanishes on $V(\mathfrak p)$ vanishes here as well.\\

\begin{exer}
	Let $X$ be a scheme and $f \in \Sh_X(X)$. Then the map
	$U \mapsto f|_U \Sh_X(U)$ for every affine open subset $U$ defines a
	sheaf of ideals on $X$, denoted $f\Sh_X$. In addition, $\Supp f\Sh_X$ is
	closed.
\end{exer}
\begin{proof}
	It is easy to see that $f\Sh_X$ inherits the uniqueness condition from
	$\Sh_X$. On the other hand, suppose that we have an open set $U$ of $X$,
	an affine open cover $\{ U_i \}$ of $U$, and a family of sections
	$f|_{U_i}s_i$ that agrees on intersections. Then we may glue these into
	a section $s$ of $\Sh_X(U)$ that restricts to $f_{U_i}s_i$ on every
	$U_i$, which we contest is contained in $f\Sh_X(U)$. Indeed, suppose
	that $s$ is not in this ideal. Then for any $a \in \Sh_X(U)$, there is
	some $x \in U$ such that $s_x - (af)_x \neq 0$, but this contradicts the
	fact that $f|_{U_i}s_i$ is in $f\Sh_X(U)$.\\

	To see that $\Supp f\Sh_X$ is closed, observe that $\Supp f\Sh_X$ is
	just the set of $x \in X$ such that $f_x \neq 0$. Since the complement
	is open, the statement holds.
\end{proof}

\begin{exer}
	Let $Y$ be a scheme such that $\Mor(X,Y) \cong \Hom(\Sh_Y(Y),\Sh_X(X))$
	for any affine scheme $X$. Then $Y$ is affine.
\end{exer}
\begin{proof}
	Regrettably, the easiest proof I could come up with was by categoryslop.
	Observe that for any affine scheme $Z$, we have $\Mor(Z,Y) \cong
	\Hom(\Sh_Y(Y),\Sh_Z(Z)) \cong \Mor(Z,\Spec(\Sh_Y(Y)))$. This extends to
	the category of all schemes in the obvious way, so $Y$ is a representing
	object of $\Mor(-,\Spec(\Sh_Y(Y))$ and is thus isomorphic to
	$\Spec(\Sh_Y(Y))$.
\end{proof}

\begin{exer}
	Let $A$ be a ring and consider the $A$-scheme $X =
	\Spec(A[T_1,\ldots,T_n]/I)$. The $A$-points of $X$ are in bijective
	correspondence with the $n$-tuples $\alpha = (\alpha_1,\ldots,\alpha_n)
	\in A^n$ such that $P(\alpha) = 0$ for all $P \in I$.
\end{exer}
\begin{proof}
	An $A$-point of $X$ is the same as a homomorphism $A[T_1,\ldots,T_n]/I
	\to A$ that restricts to the identity on $A$, so it is determined only
	by the values of the $T_i$. This homomorphism must send $T =
	(T_1,\ldots,T_n)$ to a family $\alpha = (\alpha_1,\ldots,\alpha_n)$
	satisfying the condition above. The statement is now obvious.
\end{proof}

\begin{exer}
	Let $X$ be a scheme over the field $k$. Let $\varphi: k[T_1,\ldots,T_n]
	\to \Sh_X(X)$ be a homomorphism of $k$-algebras and $f: X \to \Aff_k^n$
	the morphism induced by $\varphi$. Then for any rational point $x \in
	X(k)$, via the identification $\Aff_k^n = k^n$, we have $f(x) =
	(f_1(x),\ldots,f_n(x))$, where $f_i = \varphi(T_i)$ and $f_i(x)$ is the
	image of $f_i$ in $k(x) = k$.
\end{exer}
\begin{proof}
	The statement is obvious if you draw a picture and use the universal
	property of affine schemes:
	\[\begin{tikzcd}
						  &  X \arrow[dl,shift left]
						       \arrow[dd,"f"]
						  & & &
					       & \Sh_X(X) \ar[dl,shift right,
						  "x'"']
						  \\
		\Spec k \ar[ur,shift left,"x"] &{}&
						\ar[r, Leftrightarrow]
					       &{}&
						 k \ar[ur,shift right] \ar[dr]&
						 \\
						 & \Aff_k^n \arrow[ul]
						 & & &
					       & k[T_1,\ldots,T_n]
						  	\ar[uu,"\varphi"]
	\end{tikzcd}\]
\end{proof}

\begin{exer}
	Let $X$ be a scheme over a ring $A$. Let $f_0,\ldots,f_n \in \Sh_X(X)$
	be such that the $f_{i,x}$ generate the unit ideal of $\Sh_{X,x}$ for
	all $x \in X$. Then $X$ is the union of the $X_{f_i}$. Moreover,
	there is a morphism $f: X \to \Prs^n_A$ such that $f^{-1}(D_+(T_i)) =
	X_{f_i}$ and ${f|}_{X_{f_i}}$ is induced by the homomorphism
	$A[T_i^{-1}T_j]_j \to \Sh_X(X_{f_i})$ given by $T_i^{-1}T_j \mapsto
	f_i^{-1}f_j$.
\end{exer}
\begin{proof}
	Let $x \in X$ be any point. Since the $f_{i,x}$ generate the unit ideal,
	there is $j$ such that $f_{j,x}$ is not contained in the maximal ideal
	at $x$, so $x \in X_{f_j}$. On the other hand, we have a homomorphism
	$A[T_0,\ldots,T_n] \to \Sh_X(X)$ given by $T_i \mapsto f_i$. This gives
	rise to a homomorphism $A[T_i^{-1}T_j]_j \to \Sh_X(X)_{f_i} \cong
	\Sh_X(X_{f_i})$ for any $i$, hence a morphism $X_{f_i} \to D_+(T_i)$. We
	contest that this agrees on intersections. Indeed, observe that for any
	$i,j$, the morphisms $X_{f_i} \to D_+(T_i)$ and $X_{f_j} \to D_+(T_j)$,
	when restricted to $X_{f_if_j} = X_{f_i} \cap X_{f_j}$,
	both factor through $D_+(T_i) \cap D_+(T_j) = D_+(T_iT_j)$; it is then
	clear that they agree, so we have the desired morphism $X \to \Prs^n_A$.
\end{proof}

Note: I omitted the last part of 2.3.8 because it is the same as 2.3.7.

\begin{exer}
	Let $B$ be a graded ring. Let $f \in B$ be a non-nilpotent, homogeneous
	element of degree $0$ Then $B_f$ posseses a natural grading, and $D_+(f)
	= \Proj B_f$.
\end{exer}
\begin{proof}
	The grading is obvious. Recall that a prime ideal only localizes to a
	prime in $B_f$ if it does not contain $f$. Then $D_+(f)$ is just $\Proj
	B_f$.
\end{proof}

\begin{exer}
	Let $A$ be a ring, $X = \Prs^n_A$. Then $\Sh_X(X) = A$, and $X$ is
	affine iff $n = 0$.
\end{exer}
\begin{proof}
	Consider the covering $\{D_+(T_i)\}$ of $X$. Then a global section $f
	\in \Sh_X(X)$ must restrict to a family of sections $f_i =
	\Sh_X(D_+(T_i))$ that agree on intersections. In other words, it is
	identified with a family of functions $f_i = P_i(T_i^{-1}T_j)_j$ such
	that $P_i(T_i^{-1}T_k)_k = P_j(T_j^{-1}T_k)_k$ for all $i$, $j$. This
	can only happen if all the $P_i$ are constant, hence $\Sh_X(X) = A$.
	That $X$ is affine only when $n = 0$ follows immediately.
\end{proof}

\begin{exer}
	Let $B = \oplus_{d \geq 0}B_d$ be a graded ring. Let $e \geq 1$ be an
	integer.
	\begin{enumerate}[label=(\alph*)]
		\item Let us denote the graded ring $\oplus_{d \geq 0}B_{de}$ by
			C. Then $\Proj B \cong \Proj C$.
		\item Let us suppose that for every $d \geq 1$, $B_{ed}$ is
			generated by $B_e^d$, and that $B_e$ is of finite type
			over a ring $A$. Then $\Proj B$ is a projective scheme
			over $A$.
	\end{enumerate}
\end{exer}
\begin{proof}
	For part (a), observe that any homogeneous prime ideal of $B$ gives rise
	to one in $C$ and vice versa, and this correspondence is bijective. If
	$f$ is any non-nilpotent homogeneous element of $B$, a homogeneous prime
	ideal $\mathfrak p$ of $B$ does not contain $f$ iff it does not contain
	$f^e$, hence $\Proj B$ and $\Proj C$ have the  same support. It is easy
	to see that $C_{(f^e)}$ is isomorphic to $B_{(f)}$ by multiplying
	numerators and denominators by a suitable power of $f$, so the statement
	follows.\\

	Part (b) follows from some basic noggin use.
\end{proof}

\begin{exer}
	Let $B = A[X,Y,Z]$ be a polynomial ring. Let $B_d$ be the sub-$A$-module
	of $B$ generated by elements of the form $X^aY^bZ^c$ with $a+2b+3c=d$.
	\begin{enumerate}[label=(\alph*)]
		\item The $B_d$ induce a grading on $B$.
		\item As a module, $B_6$ is generated by $X^6$, $Y^3$, $Z^2$,
			$X^2Y^2$, $X^4Y$, $X^3Z$, and $XYZ$. Moreover, for all
			$d$, $B_{6d}$ is generated by $B_6^d$, so $\Proj B$ is
			a closed subscheme of $\Prs^6_A$.
	\end{enumerate}
\end{exer}
\begin{proof}
	The first part is obvious, and the generating set for $B_6$ is checked
	easily. Suppose that $B_{6(d-1)} = B_6^{d-1}$. Let $a$, $b$, $c$ be
	such that $X^aY^bZ^c \in B_{6d}$. If $a \geq 6$, $X^aY^bZ^c \in
	B_6B_{6(d-1)}$, so is an element of $B_6^d$. Otherwise, $2b + 3c
	\equiv -a \pmod{6}$, so there are non-negative integers $b'$, $c'$ such
	that $b' \leq b$, $c' \leq c$, and $a + 2b' + 3c' = 6$. It follows that
	$X^aY^bZ^c = (X^aY^{b'}Z^{c'})(Y^{b-b'}Z^{c-c'})$ is an element of
	$B_6B_{6(d-1)}$, and so $B_{6d}$ is generated by $B_6^d$. By Exercise
	3.11, $\Proj B \cong \Proj A[X^6,Y^3,Z^2,X^2Y^2,X^4Y,X^3Z,XYZ]$, and the
	right hand side is clearly a subscheme of $\Prs^6_A$.
\end{proof}

\begin{exer}
	A scheme is quasi-compact if and only if it is a finite union of affine
	schemes.
\end{exer}
\begin{proof}
	Use noggin.
\end{proof}

\begin{exer}
	Let $X$ be a quasi-compact scheme, $A = \Sh_X(X)$. The morphism $f: X
	\to \Spec A$ is dominant.
\end{exer}
\begin{proof}
	Let $\{U_i\}$ be a finite open covering of $X$ by affines. It is easy to
	see that $f$ is obtained by gluing the morphisms $U_i \to \Spec A$,
	which are in turn induced by the restriction maps $A \to \Sh_X(U_i)$. If
	$g$ is any non-nilpotent element of $A$, there is some $i$ such that
	${g|}_{U_i}$ is non-nilpotent, since the cover is finite. If $x$ is any
	element of $U_i$ such that ${g|}_{U_i}(x) = g(x)$ is non-zero, we see
	that $f(x) \in D(g)$ and so $f(X)$ is dense in $\Spec A$.
\end{proof}

\begin{exer}
	A scheme $X$ is locally noetherian if and only if any affine open
	subscheme of $X$ is noetherian.
\end{exer}
\begin{proof}
	The "only if" direction is clear. On the other hand, suppose that $X$
	is locally noetherian and let $U$ be an affine open subscheme. Then
	$U$ admits an open covering $\{U_i\}$ by noetherian affines, which we
	may assume finite by quasi-compactness.
\end{proof}

\begin{exer}
	Let $f: X \to Y$ be a morphism of schemes.
	\begin{enumerate}[label=(\alph*)]
		\item If $f$ is a closed immersion, it is quasi-compact.
		\item If $f$ is an open immersion and $Y$ is locally noetherian,
			then $f$ is quasi-compact.
	\end{enumerate}
	Let us suppose in what follows that $f$ is quasi-compact.
	\begin{enumerate}[label=(\alph*)]
		\setcounter{enumi}{2}
		\item Let $\mathcal I = \Ker(f^\sharp: \Sh_Y \to f_*\Sh_X)$.
			The ringed subspace $Z = V(\mathcal I)$ of $Y$ is a
			scheme.
		\item Let $j: Z \to Y$ be the closed immersion. Then we have a
			morphism $g: X \to Z$ such that $f = j \circ g$.
			Moreover, if $f$ decomposes into a morphism $g': X \to
			Z'$ followed by a closed immersion $j': Z' \to Y$, then
			$Z$ is a closed subscheme of $Z'$.
		\item The image $f(X)$ of $f$ is dense in $Z$. If $X$ is
			reduced, then so is $Z$.
	\end{enumerate}
\end{exer}
\begin{proof}
	Part (a) is clear. For part (b), just recall that any open affine
	subscheme $U$ of $Y$ is noetherian. It suffices to show part (c) for
	affine $Y$. Indeed, if $\{U_i\}$ is a covering of $Y$ by affine opens,
	we get an induced covering of $Z$ (not necessarily by affines). If each
	$U_i \cap Z$ is a scheme, then we can just paste the coverings. Now,
	assume $Y$ is affine. The quasi-compactness hypothesis implies that $X$
	is quasi-compact, so it admits a finite open covering $\{X_i\}$ by
	affines. Since the canonical map $\Sh_X(X) \to \oplus \Sh_X(X_i)$ is an
	injection, $I = \Ker(f^\sharp)(X) = \Ker(A \to \oplus \Sh_X(X_i))$. We
	contest that $I$ defines the sheaf $\Ker(f^\sharp)$. Indeed, the
	flatness of $A_g$ for any $g \in A$ gives us a commutative diagram
	\[\begin{tikzcd}
		0 \ar[r] & I \ar[d] \ar[r] & A \ar[d] \ar[r] &
		\oplus \Sh_X(X_i) \ar[d]\\
		0 \ar[r] & I_g \ar[r] & A_g \ar[r] & \oplus (\Sh_X(X_i) \otimes
		A_g)
	\end{tikzcd}\]
	where the rows are exact. Since the $X_i$ are affine, the bottom right
	homomorphism corresponds to the restriction $f^{-1}(D(f)) \to D(f)$ of
	$f$, and so $\Ker(f^\sharp)(D(f)) = I_f$. It follows that $Z \cong
	\Spec A/I$\\

	For part (d), let $g: X \to Z$ be the map (not morphism) such that
	$g(x)$ is the unique $z \in Z$ such that $f(x) = j(z)$. It is
	continuous. For any affine open $U$ of $Y$, define $g^\sharp(U \cap Z):
	\Sh_Z(U \cap Z) \to (g_*\Sh_X)(U \cap Z)$ to be the homomorphism such
	that $\Sh_Y(U) \to (f_*\Sh_X)(U) = (g_*\Sh_X)(U)$ factors through
	$g^\sharp(U \cap Z)$ by the quotient $\Sh_X(U) \to \Sh_Z(U \cap Z) = 
	\Sh_X(U)/\Ker(f^\sharp(U))$. It is easy to see that this defines a
	morphism of schemes $(g,g^\sharp):X \to Z$. If
	\begin{tikzcd} X \ar[r,"g'"] & Z' \ar[r,"j'"] & Y \end{tikzcd} is any
	other such factorization, observe that every $y \in Y$ such that
	$\Ker(f^\sharp)_y \neq \Sh_{Y,y}$ must be contained in $Z'$, since any
	such $y$ is an accumulation point of $f(X)$. It is easy to see that this
	implies that $Z = V(\Ker(g'^\sharp: \Sh_{Z'} \to g'_*\Sh_X))$.\\

	Part (e) follows from the same accumulation point argument.
\end{proof}

\begin{exer}
	Let $X$ be an affine algebraic variety over a field $k$. There is a
	projective variety $\overline{X}$ over $k$ such that $X$ is isomorphic
	to a dense open subscheme of $\overline{X}$.
\end{exer}
\begin{proof}
	Let $X = \Spec A = \Spec k[T_1,\ldots,T_n]/I$. It is easy to see that
	embeds as a closed subscheme of $D_+(T_0) \subseteq \Prs^n_k$, so let
	$J$ be the kernel of the corresponding ring homomorphism
	$k[T_0^{-1}T_1,\ldots,T_0^{-1}T_n] \to A$. We may extend $J$ to a
	homogeneous ideal $J' = k[T_0,\ldots,T_n,T_0^{-1}]J \cap
	k[T_0,\ldots,T_n]$, so we have a projective variety $V_+(J')$. We
	contest that $D_+(T_0)$ is dense in $V_+(J')$. Indeed, let $P$ be any
	homogeneous polynomial in $k[T_0,\ldots,T_n]\setminus\sqrt{J'}$. Then
	$D_+(P) \cap D_+(T_0) \cap V_+(J') = D_+(T_0P) \cap V_+(J')$, which is
	nonempty since $J'$ is not maximal (it does not contain $T_0$).
\end{proof}
